\chapter{Revisão da Literatura / Trabalhos Relacionados}
\label{sec:revisao}
\section{Estratégia de Busca da Literatura}
    Para a realização da revisão de literatura, foram conduzidas busca na platarforma 
    Google Acadêmico(Google Scholar) como objetivo de localizar trabalhos semelhantes 
    a este. De inicio foi realizada uma busca tendo em vista as palavras-chaves que se relacionam 
    com o tema, como: \textit{teste de penetração}, \textit{pentest}, \textit{ferrametas de segurança}, \textit{virtualzação}, 
    \textit{Docker}, \textit{Interface grafica educação}, \textit{LLM}.
     
    Para iniciar as pesquisa foi utilizada a seguinte String
    \textit{"Pentest" AND "Ferramentas" }, obtendo cerca de 570 resultados, apartir deste momento 
    foram separados os artigos de maior relevancia porém ainda não havia encontrado algum trabalho
    se assemelhasse integralmente com este.Portanto foi decido que a pesquisa no idioma inglês poderia levar
    a melhores resultados. Logo, a seguinte string foi pesquisada: \textit{"penetration testing" AND "GUI" }, desta vez 
    foram obtido 5.504 resultados que se mostraram muito úteis pelo seu alinhamento com a proposta deste trabalho, pórem devido a grande 
    quantidade não foi possivel ler todos. Para refinar um pouco mais a busca, adicionou-se mais um termo a string
    ficando da seguinte forma, \textit{"penetration testing" AND "GUI" AND "Education" AND "LLM" }, resultou em 118 resultados
    sobre os quais foram obtidos os demais artigos e livros relacionados.

\section{Trabalhos Correlatos}

\section{Comparação e Lacunas (Diferenciais do VAPOREON)}